% Dissertação — PPGCONT/UFRGS. Usa a mesma formatação do projeto
% (../latex/ppgcont.sty). Os capítulos 1 a 3 partem do texto de ../projeto/.
\documentclass[12pt,a4paper]{article}
% Preâmbulo da dissertação. A formatação vem de ../latex/ppgcont.sty, a mesma
% do projeto (modelo PPGCONT/UFRGS 2024). Aqui ficam só os ajustes do documento.
\usepackage{../latex/ppgcont}
\addbibresource{referencias/referencias.bib}

% Pastas de artefatos gerados pelo Python (nunca editar à mão)
\graphicspath{{figuras/}}
\newcommand{\tabelagerada}[1]{\input{tabelas/#1.tex}}

% Notação recorrente (proposta v5)
\newcommand{\PD}{\mathrm{PD}}
\newcommand{\LGD}{\mathrm{LGD}}
\newcommand{\EAD}{\mathrm{EAD}}
\newcommand{\ECL}{\mathrm{ECL}}


\titulo{Modelagem adaptativa de probabilidades de inadimplência com
  \textit{Machine Learning}: experimentos de transferência entre carteiras
  sintéticas e implicações para a mensuração de perdas esperadas de crédito}
\autor{Nome e Sobrenome}                                % TODO: nome completo
\orientador{Orientador: Prof. Dr. Nome Sobrenome}       % TODO: orientador(a)
\ano{2027}
% TODO: conferir a redação oficial da natureza do trabalho para a dissertação
\natureza{Dissertação submetida ao Programa de Pós-Graduação em Controladoria e
  Contabilidade da Faculdade de Ciências Econômicas da Universidade Federal do
  Rio Grande do Sul, como requisito parcial para a obtenção do título de Mestre
  em Controladoria e Contabilidade.}

\begin{document}
\capa
\folhaderosto
% Resumo: escrever ao final. Sem resultados até a execução dos experimentos.
\clearpage
{\centering\bfseries RESUMO\par}\vspace{\baselineskip}

\noindent Texto do resumo.

\vspace{\baselineskip}
\noindent\textbf{Palavras-chave:} perda esperada de crédito; probabilidade de
inadimplência; aprendizado de máquina; transferência de aprendizado; simulação de
Monte Carlo.

\sumario

\section{Introdução}
\label{cap:introducao}
\aescrever
% Fonte: docs/proposta/proposta-dissertacao-v5-2026-11-14.md, seção 1 e 2.

\subsection{Contextualização}
\aescrever
\subsection{Problema científico}
\aescrever
\subsection{Justificativa e relevância}
\aescrever
\subsection{Delimitação}
\aescrever
\subsection{Objetivos}
\aescrever
\subsubsection{Objetivo geral}
\aescrever
\subsubsection{Objetivos específicos}
\aescrever

\section{Referencial teórico}
\label{cap:referencial}
\aescrever
% Fonte: proposta v5, seção 3 e 5; mini RSL Q0–Q6 em docs/revisao/.

\subsection{Qualidade da informação contábil e incerteza de estimação}
\aescrever
\subsection{Transferência de aprendizado, calibração e heterogeneidade}
\aescrever
\subsection{Estudos de simulação e avaliação de métodos}
\aescrever
\subsection{Normas de perdas esperadas e enquadramento regulatório}
\aescrever
\subsection{Lacuna e hipóteses}
\aescrever

\section{Metodologia}
\label{cap:metodologia}
\aescrever
% Fonte: proposta v5, seções 4, 6, 7 e 9; protocolo experimental v5.

\subsection{Delineamento e lógica geral}
\aescrever
\subsection{Instituições virtuais e base sintética}
\aescrever
\subsection{Mecanismos geradores de dados}
\aescrever
\subsection{Estratégias de estimação (R0–R4)}
\aescrever
\subsection{Treinamento e validação temporal}
\aescrever
\subsection{Medidas de desempenho}
\aescrever
\subsection{Da PD à perda esperada sintética}
\aescrever
\subsection{Replicações e análise inferencial}
\aescrever
\subsection{Validade, robustez e reprodutibilidade}
\aescrever

\section{Resultados}
\label{cap:resultados}
\aescrever
% Preencher somente após a execução dos experimentos.
% Tabelas e figuras vêm de dissertacao/tabelas e dissertacao/figuras,
% geradas pelo pacote ecl_suficiencia.reporting. Exemplo de uso:
%   \tabelagerada{contraste-h1}
%   \includegraphics[width=\textwidth]{fig-calibracao.pdf}

\section{Discussão}
\label{cap:discussao}
\aescrever

\section{Considerações finais}
\label{cap:consideracoes}
\aescrever
\subsection{Limitações}
\aescrever
\subsection{Agenda futura}
\aescrever


\printbibliography

\appendix
\section{Parâmetros dos mecanismos geradores de dados}
\label{ap:dgp}
\aescrever
% Tabela de parâmetros e proveniência dos DGPs (gerada a partir de configs/).

\end{document}
